\documentclass[conference]{IEEEtran}
\IEEEoverridecommandlockouts
% The preceding line is only needed to identify funding in the first footnote. If that is unneeded, please comment it out.
\usepackage{cite}
\usepackage{amsmath,amssymb,amsfonts}
\usepackage{algorithmic}
\usepackage{graphicx}
\usepackage{textcomp}
\usepackage{caption}
\usepackage{xcolor}
\usepackage{booktabs}
\usepackage{tabularx}
\usepackage{cuted}
\usepackage{multirow}
\usepackage{xspace}
\usepackage{colortbl}
\usepackage{flushend}

\usepackage[caption=false,font=footnotesize]{subfig}
\def\BibTeX{{\rm B\kern-.05em{\sc i\kern-.025em b}\kern-.08em
    T\kern-.1667em\lower.7ex\hbox{E}\kern-.125emX}}

% \gap{...} marks an open item that still needs a decision, citation, or
% missing content before submission. Search this file (and methodology.tex /
% experiments.tex, which are \input below) for "\gap{" to find every one.
\providecommand{\gap}[1]{\textbf{[GAP: #1]}}
\newcommand{\method}{VICTR\xspace}
\input{macros.tex}

\begin{document}
\raggedbottom
\setlength{\textfloatsep}{4pt plus 2pt minus 2pt}
\setlength{\floatsep}{4pt plus 2pt minus 2pt}
\setlength{\intextsep}{4pt plus 2pt minus 2pt}
\setlength{\dbltextfloatsep}{4pt plus 2pt minus 2pt}
\setlength{\dblfloatsep}{4pt plus 2pt minus 2pt}

\let\OLDthebibliography\thebibliography
\renewcommand\thebibliography[1]{
  \OLDthebibliography{#1}
  \setlength{\itemsep}{1pt}
  \setlength{\parskip}{0pt}
}

\renewcommand{\dbltopfraction}{0.95}
\renewcommand{\dblfloatpagefraction}{0.8}
\renewcommand{\topfraction}{0.95}
\renewcommand{\textfraction}{0.05}
\setcounter{dbltopnumber}{2}

\title{VICTR: Value-Informed Context Retrieval for Vision-Language-Action Models\\
}
\author{\IEEEauthorblockN{Anonymous Authors}}

% \author{
% \IEEEauthorblockN{Luca Macesanu\textsuperscript{*,1}, Tianluo Zhang\textsuperscript{*,1}, Kevin Angers\textsuperscript{1}, Aditya Patil\textsuperscript{1}, Archit Sharma\textsuperscript{1}}
% \IEEEauthorblockA{\textsuperscript{1}Affiliation}
% \thanks{\textsuperscript{*}Equal contribution.}
% }

\maketitle

% \begin{strip}
%     \centering
%     \includegraphics[width=0.98\textwidth]{figures/blicl_pull.pdf}
%     \captionof{figure}{\textbf{Overview of \method (placeholder).}
%     \textbf{(Left)} At deployment, a user supplies a language instruction and a
%     small pool of teleoperated demonstrations ($\leq$10) for a new or
%     out-of-domain task. \textbf{(Center)} An In-Context Value Function
%     Estimator (IC-VFE) conditions on a full context demonstration to produce a
%     calibrated progress estimate $\hat v_t$ for the current observation; this
%     estimate is fused with visual similarity to retrieve the most relevant
%     context chunk from the demonstration pool. \textbf{(Right)} The retrieved
%     chunk conditions the IC-VLA's next action-chunk prediction, closing
%     the loop without any parameter updates. \gap{placeholder figure --
%     replace with the final architecture diagram; caption is a first draft
%     and should be revised once Section~\ref{sec:ic-vla-architecture} and
%     Section~\ref{sec:retrieval} are finalized.}}
%     \label{fig:pull-figure}
% \end{strip}

\begin{abstract}
Vision-language-action (VLA) models are increasingly capable generalists, but
adapting a pretrained VLA to a new embodiment, environment, or task still
typically requires collecting many demonstrations, mixing them with the
original training data, and fine-tuning the base policy -- an expensive loop
that must be repeated for every new domain. In-context learning (ICL) offers an alternative -- where a model's behavior is adapted at inference time by conditioning on examples supplied as part of its input, rather than by updating its weights. Prior work has shown that VLAs can be post-trained to use a
handful of demonstrations placed directly in their context window (i.e.\
prepended as extra input tokens the policy attends to when predicting its
next action, with no change to its weights), avoiding further gradient
updates altogether. A central open question for in-context
VLAs is what to retrieve into that context -- i.e., which stored
demonstration segment to select and place in the window for a given query --
and how: existing methods rely primarily on visual similarity between the
current observation and candidate demonstration frames. We argue that task-relative \emph{progress} is an
equally important and underexplored retrieval signal, and that value
function estimators (VFEs) -- widely used to provide reward signals for VLA
reinforcement learning -- can themselves be made in-context learners. We
present three contributions: (1) the first In-Context Value Function
Estimator (IC-VFE), which conditions on a full-length demonstration to
produce calibrated progress estimates on in- and out-of-domain tasks without
task-specific fine-tuning; (2) the use of estimated state value, fused with
visual similarity via a learned scalar combination, as a retrieval signal for
in-context VLA policies, extending prior work that relies on visual
similarity alone; and (3) \method, a policy that composes both signals at
inference time, together with a rigorous real-world evaluation across
\NumTasks{} tasks on a bimanual manipulator, ablating each design element
against a RICL baseline.
\end{abstract}

\begin{IEEEkeywords}
Vision-Language-Action (VLA) Models, In-Context Learning (ICL), Retrieval-Augmented Generation (RAG), Value Function Estimation, Reinforcement Learning, Bimanual Manipulation
\end{IEEEkeywords}

\section{Introduction}
\label{sec:intro}

% \paragraph{Scalable deployment challenges for VLAs.}

Vision-language-action (VLA) models are a promising 
approach to achieving generalist manipulation in the real world.
Much of their success has come from exposing them to increasingly diverse
datasets~\cite{pi05,pi06}, and reinforcement learning has recently been shown
to boost their performance on specific tasks beyond the ceiling of imitation
learning alone~\cite{pi06}. 
Yet, even with increasingly large pretraining efforts, new tasks or robot embodiments still may exist outside of the pretraining domain.
Recent distributed real-world evaluations bear this
out at scale: even strong generalist policies exhibit stark decreases in performance especially on long-horizon and out-of-distribution tasks once evaluated outside the environment or task set it was originally trained on ~\cite{roboarena}.
Simply collecting a large set of expert demonstrations in the new domain and fine-tuning the policy is slow and must be repeated for every new domain, underscoring the need for more efficient adaptation mechanisms.

% Imitation learning alone is also typically insufficient to
% reach the throughput required of a robust VLA on a specific task, motivating
% on-policy practice with occasional human intervention~\cite{pi06}. The
% central difficulty is the end-to-end cost of adapting a VLA to a new domain:
% collect many demonstrations, retrain the base policy, collect many rollouts,
% and fine-tune again. This work targets a much shorter loop: collecting few
% demonstrations ($\leq$10, rather than hundreds), performing real rollouts in
% which those demonstrations are only ever supplied to a fixed-parameter
% policy as context -- never used to update its weights -- and requiring at
% most a single fine-tuning step to reach comparable results. We propose a bi-level retrieval
% policy with an in-context value function estimator to enable this loop.

% \paragraph{In-context learning in VLAs.}

In-context learning (ICL), which is the practice of supplying examples tasks directly into the context window of a model without updating its weights, has
already proven successful in Vision-language models (VLMs) ~\cite{flamingo}.
This idea has recently been extended from VLMs to VLAs: prior work trains policies that
use a handful of demonstrations as context, boosting in- and out-of-domain
performance without any fine-tuning stage. A central open question is what
to retrieve -- i.e., which stored demonstration segment a policy should
select and place into its context window for a given query -- and how;
existing methods rely on visual similarity, alone or combined with
proprioception~\cite{ricl,mimicdroid}. We argue that task-relative progress
is an important and underexplored complement to visual similarity.

% \paragraph{Value function estimators.}
Value function estimators (VFEs) have surged in popularity for VLAs, since
they provide dense reward signals that can be used for a variety of downstream uses, including Reinforcement Learning. 
General-purpose VFEs built on VLM backbones~\cite{robometer,topreward,sarm} are an active area of work, but they still improve with task-specific fine-tuning.
We observe that these transformer-based VFEs can instead benefit from the same ICL machinery already applied to VLMs and VLAs: to the best of our knowledge, this work
proposes the first in-context value function estimator, conditioning on a
full-length demonstration to boost value-assignment quality on
out-of-domain or long-horizon tasks without task-specific fine-tuning.

% \paragraph{\method.}

This work presents \textbf{V}alue-\textbf{I}nformed \textbf{C}on\textbf{t}ext \textbf{R}etrieval for VLAs. At test time, \method accepts a language query
and the robot's current observation, and uses our
In-Context Value Function Estimator (IC-VFE) to obtain an estimated value
$\hat v_t$ for that observation. 
This estimate is fused with visual similarity, following RICL's retrieval formulation~\cite{ricl}, to retrieve a context chunk from a demonstration pool consisting of a small set of expert demonstrations (1--10). The retrieved context then conditions the VLA's next action-chunk prediction. To enable both the IC-VFE and the
IC-VLA to use demonstrations effectively, both are meta-trained on the same
held-out context pool per task, using oracle value annotations for retrieval
during training; at deployment, those oracle values are replaced by the
IC-VFE's own estimates. We evaluate the full method on \NumTaskPairs{} in-domain and \NumTaskPairs{}
out-of-domain tasks (Section~\ref{sec:demo-pools}).

In summary, this work makes three primary contributions:
\begin{enumerate}
    \item The use of estimated state value as a retrieval signal for
    in-context VLA learning, fused with visual similarity, extending prior work that relies on visual
    similarity alone.
    
    \item The first In-Context Value Function Estimator, which takes a full-length demonstration as context and provides superior value estimation on in- and out-of-domain tasks relative to current foundation
    value-estimation approaches.

    \item A rigorous evaluation of \method, an ablation of its design elements, and a comparison against state-of-the-art in-context learning methods on a real bimanual manipulator trained on \NumTrainHours{} hours
    across \NumTasks{} tasks of teleoperated data.
\end{enumerate}

\section{Related Work}
\label{sec:related-work}

\paragraph{In-context learning for embodied agents.} Prior work has explored injecting in-context adaptability into VLAs by retrieving and placing demonstration segments directly in a policy's context window~\cite{ricl}, learning single-demonstration ICL for non-VLA policies via next-token prediction~\cite{icrt}, and performing in-context imitation learning with explicit visual reasoning~\cite{iclr}. 
Other work approaches ICL from unstructured human play data~\cite{mimicdroid}, studies how retrieved action consequences shown in-context help VLAs generalize~\cite{reflectivevla}, or replaces sequence-conditioned retrieval with explicit point-correspondence matching between a single demonstration and the current scene~\cite{matchingpolicy}. 
A parallel line of work scales context implicitly rather than by retrieval, compressing long interaction histories into online-updated fast weights via test-time training~\cite{robottt}; \method instead keeps policy parameters fixed at test time and relies entirely on what is retrieved into the context window.
\method builds most directly on the retrieval-augmented formulation of RICL~\cite{ricl}, extending its visual-similarity-only retrieval signal with an in-context value estimate.

\paragraph{Value and reward modeling for VLAs.} A growing body of work trains
general-purpose value/reward models for robot manipulation, including
stage-aware reward modeling~\cite{sarm} and other VLM-backboned value
estimators~\cite{robometer,topreward}. RECAP~\cite{pi06} shows that a
distributional value function trained on episode-level success labels can
drive advantage-conditioned policy improvement at scale; our IC-VFE adopts a
similar distributional value-bin formulation
(Section~\ref{sec:ic-vfe-architecture}) but conditions the prediction on an
in-context demonstration rather than the query observation alone. Closest in
spirit is VITA~\cite{vita}, which also targets out-of-domain value
calibration but does so by taking online gradient steps on a
self-supervised loss at test time; the IC-VFE instead performs this
calibration entirely through attention over an in-context demonstration,
requiring no test-time parameter updates.
Advantage-Weighted Flow Matching~\cite{awfm} similarly uses value annotations
to bias flow-matching-based policy training toward high-reward behavior.

\paragraph{Retrieval in robot learning.} Beyond ICL specifically, retrieval
has been used to augment behavior-cloning policies with relevant
offline data~\cite{mimicgen}, and dedicated retrieval-augmented VLA
architectures have been proposed for general manipulation~\cite{retrievalvla}.
LEACL explores LLM-enhanced automatic curriculum construction for robot
learning~\cite{leacl}. \method differs from this line of work in retrieving
based on a learned, task-progress-aware signal rather than similarity alone.

% methodology.tex
% Methodology section for VICTOR paper.
% Intended to be \input{} from the main paper .tex file.



\section{Methodology}
\label{sec:methodology}

\begin{figure*}[t]
    \centering
    \includegraphics[width=0.98\textwidth]{figures/blicl_method.pdf}
    \caption{\textbf{\method method overview (placeholder).} At deployment,
    the demonstration pool $\mathcal{D}_\ell$ feeds two in-context
    estimators. The IC-VFE ($\phi$) conditions on one context demonstration
    to produce a calibrated progress estimate $\hat v_t$ for the current
    observation $o_t$; $\hat v_t$, fused with visual similarity, drives
    retrieval of a context chunk $c_t$ from the same pool. The IC-VLA
    ($\theta$) conditions on $c_t$ to predict the next action chunk,
    closing the loop without any parameter updates. \gap{placeholder
    figure -- replace with final diagram once the layout in
    Section~\ref{sec:context-retrieval} is finalized.}}
    \label{fig:method}
\end{figure*}

In this section we present an overview of \method: the problem setting, how
per-task context is constructed and retrieved, the architecture and
meta-training objective of the In-Context Value Function Estimator (IC-VFE),
and the architecture and meta-training objective of the IC-VLA policy.

\subsection{Problem Formulation}
\label{sec:problem-formulation}

We consider a multi-task vision-language-action setting. A task is specified by a
language instruction $\ell \in \mathcal{L}$, and an observation
$o_t = [X_t^1, \dots, X_t^n, q_t]$ consists of $n$ camera images and proprioceptive
state $q_t$. In our real-robot deployment $n=3$ and $q_t$ is an 18-dimensional
per-arm end-effector position and orientation, the orientation given as a 6D
rotation representation rather than a unit quaternion to avoid the antipodal
(double-cover) ambiguity a quaternion state introduces; the action chunk
$a_{t:t+H}$ is correspondingly 20-dimensional per step (per-arm delta
end-effector position relative to the chunk's first frame, absolute 6D
rotation, and a gripper command). A base
VLA policy $\pi_\theta(a_{t:t+H} \mid o_t, \ell)$ maps an observation and
instruction to an action chunk $a_{t:t+H}$.
Each demonstration $\tau = (o_0, a_0, \dots, o_T)$ is annotated offline
(Section~\ref{sec:dataset}) with a scalar progress
value $v_t \in [0, 1]$, with $v_0 \approx 0$ and $v_T = 1$ on success.

\paragraph{Demonstration pool.} Each meta-training task $\ell$ has an associated
demonstration set $\mathcal{D}_\ell$: individual $(o_t, a_{t:t+H})$ pairs from
$\mathcal{D}_\ell$ supply queries for the base flow-matching loss, and the same
per-task episode set is flattened into the demonstration chunk dictionary
$\mathcal{D}_\ell^{\text{ctx}}$ that both in-context estimators condition on
(Sections~\ref{sec:ic-vfe-objective}, \ref{sec:ic-vla-objective}) and that supplies
the retrieval pool defined in Section~\ref{sec:context-retrieval}. Unlike an
idealized held-out split, our current real-robot training arms draw
$\mathcal{D}_\ell^{\text{ctx}}$ from the same episodes used for training queries
rather than a disjoint subset; enforcing a strict split is future work. At
deployment there is no train/context distinction to begin with: the user supplies
a single small demonstration set $\mathcal{D}_\ell$ ($N \le 10$) for the new task,
and this entire set plays the role $\mathcal{D}_\ell^{\text{ctx}}$ played during
meta-training.

\paragraph{Bi-level structure.} We call the overall method \method\ because two
independent in-context estimators compose at inference time, both conditioned on
the same demonstration pool but in different forms (Section~\ref{sec:context-retrieval}):
(1) the IC-VFE conditions on one full context demonstration to produce a calibrated
progress estimate $\hat{v}_t \in [0,1]$ for the current observation, and (2) that
estimate, fused with visual similarity, drives a retrieval step that selects a context
chunk $c_t$ consumed by the IC-VLA to predict the next action chunk. Both estimators
are meta-trained on oracle progress values $v_t$; at deployment, oracle values are
replaced by $\hat{v}_t$ from the trained IC-VFE, so downstream retrieval quality is
bounded by IC-VFE accuracy.

\subsection{Preliminaries}
\label{sec:preliminaries}

Both the base VLA and the IC-VLA follow the $\pi_{0.5}$ recipe of a VLM backbone
paired with a flow-matching action expert
that predicts the continuous action chunk $a_{t:t+H}$~\cite{pi05}. Given a denoising
timestep $\eta \in [0,1]$, a noised action chunk $a_{t:t+H}^\eta = \eta\, a_{t:t+H} +
(1-\eta)\,\epsilon$ with $\epsilon \sim \mathcal{N}(0, I)$, and a velocity field
$u_\theta$ predicted by the action expert conditioned on $(o_t, \ell)$ (and, for the
IC-VLA, the retrieved context $c_t$), the flow-matching loss is
\begin{equation}
    L_{\text{flow}}\big(a_{t:t+H} \mid \cdot\,; \theta\big)
    = \mathbb{E}_{\eta, \epsilon}\Big[\big\|u_\theta(a_{t:t+H}^\eta, \eta \mid \cdot) -
    (a_{t:t+H} - \epsilon)\big\|^2\Big] ,
    \label{eq:flow-loss}
\end{equation}
which we reuse, without modification, as the action-chunk loss term for both the base
VLA fine-tune and the IC-VLA (Section~\ref{sec:ic-vla-objective}).

\subsection{Context Construction and Retrieval}
\label{sec:context-retrieval}

The demonstration pool $\mathcal{D}_\ell^{\text{ctx}}$ (or $\mathcal{D}_\ell$ at
deployment; Section~\ref{sec:problem-formulation}) feeds the two estimators in two
different forms. First, one whole demonstration $\tau_{\text{ctx}} \sim
\mathcal{D}_\ell^{\text{ctx}}$ is sampled and given to the IC-VFE as the context
sequence it conditions on to produce $\hat v_t$ (Section~\ref{sec:ic-vfe-architecture}).
Second, the pool is broken apart into a frame-level retrieval buffer from which the
IC-VLA's retrieval step selects a context chunk $c_t$. This mirrors the ``priming
vs.\ retrieval buffer'' split in RICL~\cite{ricl}, except here the same held-out set
serves both roles for both estimators, keeping the two meta-training procedures
aligned on identical context data.

\paragraph{Chunk dictionary construction.} We flatten the $N$-demonstration pool into
a \emph{demonstration chunk dictionary}: a key--value store whose keys are
DINO-v2~\cite{dinov2} embeddings of a chunk's first frame, and whose values are the
chunk contents (per-frame image, proprioception, and action). To
build the dictionary, each of the pool's demonstrations is split into overlapping
chunks of fixed length $L$ (in frames, $L=10$ in our implementation) with a
50\%-overlap stride $S = L/2$ between consecutive chunks; a demonstration of $T$
frames contributes $\lfloor (T-L)/S \rfloor + 1$ chunks, and each chunk carries its
oracle value (meta-training) or IC-VFE-estimated value (deployment). Because
pool chunks overlap 50\% within an episode, temporally-adjacent chunks often have
near-identical DINO-v2 embeddings; to prevent unconstrained top-$k$ retrieval from
collapsing onto several near-duplicate chunks from the same episode, retrieval
enforces an episode-diversity constraint, selecting at most one chunk per source
episode among the top-$k$ neighbors.

\begin{figure}[t]
    \centering
    \includegraphics[width=\linewidth]{figures/demo_dictionary.pdf}
    \caption{\textbf{Demonstration chunk dictionary construction (placeholder).}
    Illustration of splitting a demonstration into overlapping chunks of size $L$ and
    indexing each by the DINO embedding of its first frame. \gap{placeholder figure --
    replace with final diagram.}}
    \label{fig:chunk-dict}
\end{figure}

\paragraph{Retrieval.} A retrieval function
\begin{equation}
    f_{\text{retrieve}}\big(o_t, \hat{v}_t, \mathcal{P}_\ell\big) \;\rightarrow\;
    c_t = \big(n^{(1)}, \dots, n^{(k)}\big)
    \label{eq:retrieval-fn}
\end{equation}
queries the chunk dictionary $\mathcal{P}_\ell$ and selects the $k$ nearest neighbor
chunks $n^{(i)} = (o', a', v')$ (subject to the episode-diversity constraint above),
ranked by a fused score over visual similarity $d_{\text{vis}}(o_t, o')$ (DINO-v2
embedding $\ell_2$ distance, as in RICL) and value alignment $|\hat{v}_t - v'|$. Our
vision+value \method variant combines the two by independently $z$-scoring
$d_{\text{vis}}$ and $|\hat{v}_t - v'|$ across the current candidate pool -- so the
two signals, which live on very different scales, contribute comparably -- and
summing the result, $g(d_{\text{vis}}, |\hat v_t - v'|) = z(d_{\text{vis}}) +
z(|\hat v_t - v'|)$, rather than a fixed weighting such as RICL's distance-decay
$e^{-\lambda d}$. This fixed fusion is a simpler alternative to a jointly-trained
scalar fusion $g_\psi$ -- a small 2-layer MLP over $[d_{\text{vis}}, |\hat v_t -
v'|]$, trained end-to-end with the IC-VLA via a soft, temperature-annealed
relaxation of top-$k$ selection since hard top-$k$ is not itself differentiable --
which we leave to future work (Section~\ref{sec:limitations}).

\paragraph{Context window token layout.} The $k$ retrieved neighbor chunks $c_t$ are
placed in the IC-VLA's context ordered farthest-to-nearest, so that the nearest
neighbor immediately precedes the query observation tokens. Each neighbor chunk
contributes one image per input camera (all taken from the chunk's first frame)
together with a single combined text/action token span, produced by the base
$\pi_{0.5}$ tokenizer applied to the neighbor's own task string, its initial
proprioceptive state (digitized directly into the prompt text -- the same mechanism
$\pi_{0.5}$ already uses to embed the query's own state), and its
action-horizon-length action window (FAST-tokenized); this reuses $\pi_{0.5}$'s
existing text/action embedding path exactly, introducing no new token type. Each
neighbor block receives a learned neighbor-rank embedding (added to every token in
the block) so the IC-VLA can in principle distinguish nearer from farther context;
all real-robot runs reported in this work use $k=1$ (Section~\ref{sec:limitations}),
so this embedding currently only marks the presence of retrieved context rather than
distinguishing among multiple ranks. This layout gives the context-conditioned
policy
\begin{equation}
    \pi_\theta\big(a_{t:t+H} \mid o_t, \ell, c_t\big) .
    \label{eq:ic-vla-policy}
\end{equation}

\subsection{IC-VFE Architecture and Demonstration Bin Tokenization}
\label{sec:ic-vfe-architecture}

Rather than estimating value from the query observation alone (as in RECAP's
$p_\phi(V \mid o_t, \ell)$~\cite{pi06}), the IC-VFE additionally conditions on one full
context demonstration's downsampled frame sequence $\tau_{\text{ctx}}$
(Section~\ref{sec:context-retrieval}):
\begin{equation}
    p_\phi\big(V \mid o_t, \ell, \tau_{\text{ctx}}\big) \in \Delta^B ,
    \label{eq:ic-vfe}
\end{equation}
a distribution over $B = 201$ discretized value bins (following RECAP's distributional
value head), from which we recover a scalar estimate $\hat{v}_t = \sum_b p_\phi(V = b
\mid \cdot)\, \nu(b)$. Intuitively, the context demonstration lets the IC-VFE
calibrate what progress looks like for this specific task, object, or embodiment
instance before scoring the live observation, rather than relying purely on
generalization from pretraining. A no-context ablation removes only the demonstration
sequence while retaining the same current-observation encoder and prediction head.

\paragraph{Architecture.} The context demonstration is downsampled end-anchored (most
recent frames preferred, at most 200 slots) into a fixed number of temporal
\emph{context slots}, each visually tokenized by a frozen Gemma 3 vision tower and a
trainable patch merger (compressing 256 pooled visual tokens per camera to 16, and
projecting into the 640-dimensional Gemma 270M embedding space), alongside a
proprioceptive token; with 3 cameras this gives 49 tokens per slot and up to 9{,}800
context tokens total. The current observation is separately encoded at two
resolutions: a trainable detail connector projects all 256 pooled tokens per camera
independently (preserving fine visual detail), while the same patch merger used for
the context path also produces a 16-token summary (placing the current frame in the
same compressed visual space as the demonstration). Context tokens, the dual-resolution
current-observation tokens, current proprioceptive state, the tokenized instruction
(up to 32 tokens), and a learnable value-query token are concatenated and processed by
a Gemma 3 270M causal decoder; because the value-query token is last, its final hidden
state attends to all preceding tokens and is mapped by a value head to a categorical
distribution over the 201 value bins.

\paragraph{Demonstration-derived value-bin tokenization.} The value target is defined
by the relative progress of the current query frame within its successful
demonstration. For a successful trajectory $\tau = (o_0, \dots, o_{T-1})$, the
continuous value assigned to frame $t$ is
\begin{equation}
    y_t = \frac{t}{T-1}, \qquad T > 1 ,
\end{equation}
so the initial frame has value $0$, the terminal successful frame has value $1$, and
intermediate frames are labeled by linear relative progress; for the degenerate case
$T=1$, the single terminal frame is assigned $y_0 = 1$. We discretize this continuous
target into $B = 201$ uniformly spaced value bins. The scalar represented by bin $b
\in \{0, \dots, B-1\}$ is $\nu_b = b / (B-1)$, so adjacent bins are separated by
$1/200 = 0.005$. The target bin index for value $y_t$ is
\begin{equation}
    b_t = \operatorname{clip}\left(
        \operatorname{round}\left[y_t (B-1)\right],\, 0,\, B-1
    \right) .
\end{equation}
We refer to this operation as demonstration-derived value-bin tokenization because the
categorical target is obtained from the query frame's position in a successful
demonstration; these 201 bins are supervision classes for the value head, not tokens
appended to the multimodal input sequence.

\subsection{IC-VFE Meta-Learning Training Objective}
\label{sec:ic-vfe-objective}

Let $z_\theta(o_t, \ell, \tau_{\text{ctx}}) \in \mathbb{R}^{B}$ denote the logits
produced by the value head, where $B = 201$. The predicted categorical value
distribution is
\begin{equation}
    p_\theta(b \mid o_t, \ell, \tau_{\text{ctx}})
    = \frac{\exp z_{\theta, b}}{\sum_{b'=0}^{B-1} \exp z_{\theta, b'}} .
\end{equation}
IC-VFE is trained with categorical cross-entropy against the demonstration-derived
target bin $b_t$,
\begin{equation}
    \mathcal{L}_{\mathrm{IC\text{-}VFE}}(\theta)
    = -\mathbb{E}_{(o_t, \ell, \tau_{\text{ctx}}) \sim \mathcal{B}}
    \left[\log p_\theta(b_t \mid o_t, \ell, \tau_{\text{ctx}})\right] ,
\end{equation}
where $\mathcal{B}$ denotes the training distribution over query frames, instructions,
and same-task context demonstrations sampled from $\mathcal{D}_\ell^{\text{ctx}}$
(Section~\ref{sec:problem-formulation}). For the no-context ablation, the objective is
unchanged and the conditioning variable $\tau_{\text{ctx}}$ is omitted. The frozen
vision tower and fixed average-pooling operator are excluded from optimization;
gradients update the detail connector, shared patch merger, proprio encoder, learned
source embeddings, Gemma 270M backbone, value-query token, and value head.

At inference time, a scalar value is recovered as the expectation of the predicted
categorical distribution,
\begin{equation}
    \hat{y}_t = \mathbb{E}_{p_\theta}[\nu_b] = \sum_{b=0}^{B-1} p_\theta(b \mid o_t, \ell, \tau_{\text{ctx}})\, \nu_b .
\end{equation}
This distributional formulation preserves uncertainty over progress while providing a
scalar estimate in $[0, 1]$ when required. Mean-squared error and mean-absolute error
between $\hat{y}_t$ and $y_t$ may be reported as diagnostics, but they are not
auxiliary optimization terms: the training loss is the categorical cross-entropy
alone.

\subsection{IC-VLA Architecture and Context Chunk Tokenization}
\label{sec:ic-vla-architecture}

The IC-VLA shares its base architecture and fine-tuning regime with $\pi_{0.5}$
(Table~\ref{tab:ft-regime}), extended to accept the retrieved context chunk $c_t$
alongside the query observation, following the general recipe of placing retrieved
neighbors in context ordered by rank (as in RICL~\cite{ricl}). The
token-level layout of this context window --- how each retrieved neighbor is
tokenized, how many neighbors are retrieved, and how neighbor tokens are ordered
relative to the query --- is given in Section~\ref{sec:context-retrieval}.

\subsection{IC-VLA Meta-Learning Training Objective}
\label{sec:ic-vla-objective}

The IC-VLA is trained with the context-conditioned policy $\pi_\theta(a_{t:t+H}
\mid o_t, \ell, c_t)$ (Equation~\ref{eq:ic-vla-policy}): the action expert predicts
the continuous action chunk $a_{t:t+H}$ conditioned on the query observation,
instruction, and retrieved context, trained with the flow-matching loss
$L_{\text{flow}}$ (Equation~\ref{eq:flow-loss}), following the fine-tuning regime in
Table~\ref{tab:ft-regime} (Appendix~\ref{app:dataset-prep}) and mirroring the
stop-gradient interface between the VLM backbone and the action expert used for the
base $\pi_{0.5}$ fine-tune (Table~\ref{tab:ft-regime}). The training objective is
\begin{equation}
    \mathcal{L}_{\mathrm{IC\text{-}VLA}}(\theta)
    = \mathbb{E}_{(o_t, \ell, c_t) \sim \mathcal{B}}
    \Big[L_{\text{flow}}\big(a_{t:t+H} \mid o_t, \ell, c_t; \theta\big)\Big] ,
    \label{eq:ic-vla-loss}
\end{equation}
During
meta-training, the context chunk $c_t$ is sampled by $f_{\text{retrieve}}$
(Equation~\ref{eq:retrieval-fn}) using oracle progress values $v_t$ from
$\mathcal{D}_\ell^{\text{ctx}}$ (Section~\ref{sec:problem-formulation}); at
deployment, oracle values are replaced by the trained IC-VFE's estimate $\hat{v}_t$,
so retrieval quality at inference time is bounded by IC-VFE accuracy. \gap{confirm
accuracy.}



% experiments.tex
% Experiments and Results section for VICTOR paper.
% Intended to be \input{} from the main paper .tex file, after methodology.tex.

% \gap{...} marks a remaining open item that still needs a decision or a
% citation before submission. Search for "\gap{" to find every open item.
% Defined with \providecommand so this file is safe to \input from a main
% paper .tex that may also define \gap in its own preamble.
\providecommand{\gap}[1]{\textbf{[GAP: #1]}}

\section{Experiments and Results}
\label{sec:experiments}

We evaluate \method on \NumTasks{} meta-training tasks (\NumEpisodes{}
episodes, \NumTrainHours{} hours of teleoperated bimanual data) and a
held-out suite of \NumTaskPairs{} in-domain/out-of-domain task pairs
(\NumEvalTasks{} tasks total) on the YOR bimanual manipulator~\cite{yor}, in
five stages: (1) value-function and retrieval quality, (2) the main policy
comparison, (3) a demonstration-count ablation, (4) a retrieval--policy
mismatch and context ablation, and (5) a context-sensitivity analysis.

\subsection{Value Function and Retrieval Quality}
\label{sec:vfe-retrieval-quality}

\begin{table}[t]
    \centering
    \caption{\textbf{Value-function and retrieval quality (placeholder).}
    Value correlation is measured against human progress annotations on the
    \NumEvalTasks{} evaluation tasks; retrieval MAE uses each method's value
    predictions in place of the IC-VFE to drive retrieval
    (Section~\ref{sec:vfe-retrieval-quality}).}
    \label{tab:vfe-retrieval-results}
    \footnotesize
    \begin{tabular}{lrrrr}
        \toprule
        & \multicolumn{2}{c}{Value correlation} & \multicolumn{2}{c}{Retrieval MAE} \\
        Method & Kendall $\tau_b$ & Pearson $r$ & Position & Orientation \\
        \midrule
        Robometer (zero-shot) & X.XX & X.XX & X.XX & X.XX \\
        TOPReward             & X.XX & X.XX & X.XX & X.XX \\
        GVL                   & X.XX & X.XX & X.XX & X.XX \\
        RECAP (fine-tuned)    & X.XX & X.XX & X.XX & X.XX \\
        Robometer (fine-tuned) & X.XX & X.XX & X.XX & X.XX \\
        IC-VFE (ours)         & \textbf{X.XX} & \textbf{X.XX} & \textbf{X.XX} & \textbf{X.XX} \\
        \midrule
        RoboDopamine\textsuperscript{$\dagger$} & X.XX & X.XX & X.XX & X.XX \\
        \bottomrule
    \end{tabular}
    \\[2pt]
    \raggedright\footnotesize\textsuperscript{$\dagger$}The annotation model itself, fine-tuned as a value estimator; a sanity check on how well its own training target can be fit, not a compared baseline.
\end{table}

We compare zero-shot value function estimators (Robometer~\cite{robometer},
TOPReward~\cite{topreward}, GVL~\cite{gvl}) against estimators fine-tuned on
our full meta-training dataset with dense RoboDopamine annotations
(RECAP~\cite{pi06}, Robometer, and our IC-VFE). Each method is scored by (a)
Kendall's $\tau_b$ and Pearson correlation against human-annotated progress
values, and (b) our retrieval-quality metric: using a method's value
predictions in place of the IC-VFE to retrieve a context chunk $c=(o',a')$
via $f_{\text{retrieve}}$ (Section~\ref{sec:context-retrieval}), then
computing the MAE between $a'$ and the ground-truth action paired with the
query observation, separately for position and orientation. Both metrics for
all methods appear in Table~\ref{tab:vfe-retrieval-results}.

\gap{[Analysis]}

\subsection{Main Evaluation}
\label{sec:main-eval}

\begin{figure}[t]
    \centering
    \includegraphics[width=0.95\linewidth]{figures/main_sr_barchart.pdf}
    \caption{\textbf{Total success rate by method (placeholder).} Bars show
    in-domain and out-of-domain success rate (\%) for each compared method,
    over \NumRolloutsPerTask{} rollouts per task across all \NumEvalTasks{}
    tasks. \gap{placeholder figure -- replace with real results.}}
    \label{fig:main-sr}
\end{figure}

We compare $\pi_{0.5}$ trained for 40k and 50k steps (no retrieval),
RICL-$\pi_{0.5}$~\cite{ricl}, and three \method variants: value-based
retrieval, vision-similarity retrieval, and vision+value retrieval (our main
method). All methods are trained on the same meta-training data and
evaluated over \NumRolloutsPerTask{} rollouts per task on all \NumEvalTasks{}
tasks. We report total and partial (subtask-level) success rate, split by
in-domain and out-of-domain task, in Figure~\ref{fig:main-sr}.

\gap{[Analysis]}

\subsection{Demonstration-Count Ablation}
\label{sec:demo-count-ablation}

\begin{figure}[t]
    \centering
    \includegraphics[width=0.95\linewidth]{figures/demo_count_sr.pdf}
    \caption{\textbf{Success rate vs.\ number of context demonstrations
    (placeholder).} $x$-axis $\in \{1,3,5,7,10\}$ demonstrations, $y$-axis
    success rate (\%), pooled over 6 tasks. \gap{placeholder figure --
    replace with real results.}}
    \label{fig:demo-count}
\end{figure}

Taking the best policy from Section~\ref{sec:main-eval}, we vary the number
of demonstrations available in the retrieval buffer over $\{1, 3, 5, 7, 10\}$,
randomly selected from the available pool, on 6 long-horizon tasks (3
in-domain and their 3 out-of-domain counterparts), with 10 rollouts per
configuration ($6 \times 10 \times 5 = 300$ rollouts total).

\gap{[Analysis]}

\subsection{Retrieval--Policy Mismatch and Context Ablation}
\label{sec:mismatch-ablation}

\begin{table}[t]
    \centering
    \caption{\textbf{Train/test retrieval mismatch (placeholder).} Success
    rate (\%) when the retrieval strategy at test time differs from the one
    used during meta-training.}
    \label{tab:retrieval-mismatch}
    \footnotesize
    \begin{tabular}{lrr}
        \toprule
        Train $\backslash$ Test & Value retrieval & Vision retrieval \\
        \midrule
        Value retrieval  & XX.X & XX.X \\
        Vision retrieval & XX.X & XX.X \\
        \bottomrule
    \end{tabular}
    \quad
    \begin{tabular}{lrr}
        \toprule
        Method & Default & Swapped \\
        \midrule
        $\pi_{0.5}$ (+ context) & XX.X & XX.X \\
        \method (no context)    & XX.X & XX.X \\
        \bottomrule
    \end{tabular}
\end{table}

We measure how much retrieval choices actually matter at test time in two
ways. First, we cross the training-time and test-time retrieval strategy
(value vs.\ vision) for our two single-signal \method variants in a
$2\times2$ design. Second, we swap context presence: $\pi_{0.5}$ (trained
without context) is given a retrieved context chunk at test time, and
\method (trained with context) is deployed with no context at test time.
Both results are reported as success rate in Table~\ref{tab:retrieval-mismatch}.

\gap{[Analysis]}

\subsection{Context Sensitivity}
\label{sec:context-sensitivity}

\begin{table}[t]
    \centering
    \caption{\textbf{Context sensitivity (placeholder).} MAE between the
    action predicted from the true retrieved context and the action
    predicted from random context chunks; higher indicates a policy relies
    more on its retrieved context.}
    \label{tab:context-sensitivity}
    \footnotesize
    \begin{tabular}{lrr}
        \toprule
        Method & Position MAE & Orientation MAE \\
        \midrule
        RICL-$\pi_{0.5}$        & X.XX & X.XX \\
        Ours (value retrieval)  & X.XX & X.XX \\
        Ours (vision retrieval) & X.XX & X.XX \\
        Ours (vision + value)   & X.XX & X.XX \\
        \bottomrule
    \end{tabular}
\end{table}

For each policy, we take an observation, retrieve its context chunk via that
policy's own retrieval strategy, and record the predicted action $a_c$; we
then replace the context with randomly sampled chunks from the demonstration
pool and record the resulting action $a_{\text{rand}}$. We report the MAE
between $a_c$ and $a_{\text{rand}}$, separately for position and
orientation, in Table~\ref{tab:context-sensitivity}.

\gap{[Analysis]}




\section{Limitations and Future Work}
\label{sec:limitations}

\method's evaluation protocol has several known limitations. Success and
failure are judged against a fixed per-task rubric without inter-rater
agreement statistics, and no multiple-seed reruns are performed per
configuration (Appendix~\ref{app:eval-metrics}). Composition and confounding
generalization are currently evaluated only on the 3 long-horizon and 9
same-skill task pairs in Table~\ref{tab:task-pairs}, respectively; broader
coverage would strengthen the generalization claims in
Section~\ref{sec:intro}.

The bi-level design also introduces an unmeasured compounding-error risk: an
inaccurate IC-VFE estimate $\hat v_t$ can steer retrieval toward a
mismatched context chunk, degrading policy performance independent of the
policy's own quality. Section~\ref{sec:vfe-results} reports IC-VFE accuracy
in isolation; an ablation measuring downstream policy sensitivity to VFE
error would clarify this failure mode.

Future work includes supervised fine-tuning on top of RL100-style data
collection, training the retrieval fusion $g_\psi$
(Section~\ref{sec:context-retrieval}) jointly with the IC-VLA rather than the
fixed fusion used by all reported real-robot runs, extending \method's
context beyond the single retrieved chunk ($k=1$) those same runs use -- a
memory-bound constraint of the current implementation, not a limit of the
formalism in Section~\ref{sec:context-retrieval} -- toward multi-neighbor and
longer-horizon composed retrieval, and evaluating \method's robustness on a
common simulation benchmark such as CALVIN as a further check beyond the
real-world evaluation. \method is currently training-free at test
time by design, relying entirely on retrieval into a fixed-parameter
context window; recent theoretical and empirical work shows that pairing
in-context learning with lightweight test-time training can further reduce
the number of demonstrations needed to reach a target accuracy, both in
general transformer ICL settings~\cite{ttt-icl-provably,ttt-fewshot} and,
concretely, in VLA policies via fast-weight context
scaling~\cite{robottt} and test-time-trained visual
foresight~\cite{ttt-visual-foresight}. A natural extension of \method is to
let the retrieved context also drive a small number of test-time gradient
steps on the IC-VFE or the policy itself, rather than treating retrieval and
inference as fully separate stages.

\section{Conclusion}
\label{sec:conclusion}

\gap{write after Results and Analysis are finalized with real numbers.} This
work presents \method, a bi-level in-context learning method that pairs an
In-Context Value Function Estimator with value- and vision-aware retrieval to
adapt a pretrained VLA to new tasks from as few as one to ten demonstrations,
without task-specific fine-tuning. We evaluate \method against a RICL
baseline on a real bimanual manipulator across in-domain,
novel-object, and long-horizon composition tasks.

\section*{Acknowledgment}
\gap{fill in acknowledgments once finalized.}

\bibliographystyle{IEEEtran}
\bibliography{references}
% \gap{references.bib does not yet exist in this directory and needs entries
% for at least the following keys used across this file, methodology.tex, and
% experiments.tex: pi05, pi06, ricl, dinov2, icrt, iclr, grrl, mimicdroid,
% mimicgen, retrievalvla, reflectivevla, sarm, topreward, robometer, awfm,
% leacl, progvla, libero, yor.}

% \clearpage
% \input{appendix.tex}

\end{document}
 
